\subsection{Reranking}

Il risultato della ricerca ibrida è però, tipicamente, solo una prima scrematura. Per restare efficienti su un
indice di grandi dimensioni, entrambe le sue componenti valutano ciascun frammento senza che un modello lo legga
mai insieme alla domanda. BM25 si limita a contare le parole in comune; nella componente densa, l'embedding di
ciascun frammento viene calcolato una volta per tutte, in fase di indicizzazione e quindi prima ancora di conoscere
la domanda, e al momento della ricerca viene soltanto confrontato con l'embedding della domanda tramite la
similarità coseno (Sottosezione~\ref{subsec:embedding}). Il vettore di un frammento deve perciò riassumerne l'intero
contenuto senza sapere quale aspetto interesserà a chi interroga il sistema: un frammento sulle tasse universitarie
è rappresentato sempre dallo stesso vettore, sia che la domanda riguardi la scadenza della seconda rata, sia che
riguardi l'importo previsto per una specifica fascia di reddito. Questa compressione comporta inevitabilmente una
perdita di dettaglio, e un frammento solo genericamente pertinente può finire in cima ai risultati davanti a quello
che contiene davvero la risposta. Un secondo stadio, il \emph{reranking}, corregge questo limite impiegando un
\emph{cross-encoder}: un modello che riceve domanda e documento candidato concatenati in un unico input e ne stima
direttamente la pertinenza reciproca, a un costo computazionale più alto ma applicato solo al pool già ristretto dei
migliori candidati della prima fase. Nogueira e Cho \cite{nogueira2019bert} sono tra i primi a mostrare, adattando
BERT a questo compito, che un cross-encoder di questo tipo migliora in modo sostanziale la qualità del ranking finale
rispetto alla sola similarità vettoriale.

Il cross-encoder rappresenta un estremo: domanda e documento vengono elaborati insieme, e ogni token dell'una
può essere messo in relazione con ogni token dell'altro in tutti gli strati della rete, al prezzo di dover
ripetere l'intero calcolo per ogni coppia. Esistono però soluzioni intermedie, che ottengono un confronto più dettagliato di
quello tra due singoli vettori senza sostenere il costo di un cross-encoder. Khattab e Zaharia \cite{khattab2020colbert} propongono
con ColBERT un'alternativa intermedia, la \emph{late interaction}: invece di comprimere ogni testo in un unico
vettore, ogni token di domanda e documento mantiene il proprio embedding, e la similarità complessiva si ottiene
aggregando le similarità massime tra i singoli token delle due sequenze --- un calcolo più ricco della sola
similarità tra due vettori aggregati, ma più leggero di un cross-encoder pieno, perché gli embedding dei documenti
possono comunque essere pre-calcolati e indicizzati in anticipo, mentre un cross-encoder richiede di rielaborare
congiuntamente domanda e documento a ogni confronto.